\section{The Diamond Principle}

We now introduce a principle that guesses every subset of $\lambda$ correctly on a stationary set of stages. It is what we will use to build a Suslin tree in \Cref{Ch4:Sec:Suslin_Trees}.

\begin{boxdefinition}[$\Diamond$]
    Let $\lambda$ be a regular, uncountable cardinal. Let $E \ssq \lambda$. We define \textbf{$\Diamond_{\lambda}(E)$} to be the statement that there is some sequence $\grpres{A_{\alpha}}{\alpha \in E}$ such that the following both hold:
    \begin{enumerate}
        \item Each $A_{\alpha}$ is a subset of $\alpha$
        \item For every $X \ssq \lambda$, the set $S_X = \setst{\alpha \in E}{\alpha \cap X = A_{\alpha}}$ is a stationary subset of $\lambda$.
    \end{enumerate}
    $\Diamond_{\lambda}(E)$ is shorthand for the statement that there is a $\Diamond_{\lambda}(E)$-sequence, and $\Diamond_{\lambda}$ is shorthand for the statement that there is a $\Diamond_{\lambda}(\lambda)$-sequence.
\end{boxdefinition}

There are some standard abbreviations.

\begin{boxnotation}\hfill
    \begin{enumerate}
        \item $\Diamond_{\lambda} \equiv \Diamond_{\lambda}(\lambda)$
        \item $\Diamond(E) \equiv  \Diamond_{\omega_1}(E)$
        \item $\Diamond \equiv \Diamond_{\omega_1}\of{\omega_1} = \Diamond_{\omega_1}$
    \end{enumerate}
\end{boxnotation}

A $\Diamond_{\lambda}$-sequence guesses every bounded set stationarily often.

\begin{boxlemma}
    Suppose $X$ is a bounded subset of $\lambda$ (so we have $\beta < \lambda$ such that $X \ssq \beta$). Assume $\grpres{A_{\alpha}}{\alpha < \lambda}$ is a $\Diamond_{\lambda}$-sequence. Then $\setst{\alpha < \lambda}{A_{\alpha} = X}$ is stationary in $\lambda$.
\end{boxlemma}

This already says something about cardinal arithmetic.

\begin{boxcorollary}
    $\Diamond_{\lambda} \to 2^{<\lambda} = \lambda$.
\end{boxcorollary}

In particular, when $\lambda = \omega_1$, we have $2^{< \omega_1} = 2^{\omega} = \omega_1$, so CH holds. In other words, diamond implies CH.

The sets $S_X$ themselves are nearly disjoint.

\begin{boxlemma}
    Assume $\Diamond_{\lambda}(E)$. Then $\grpres{S_X}{X \ssq \lambda}$ (where $S_X$ is defined from the witness to $\Diamond_{\lambda}(E)$) splits $E$ into $2^{\lambda}$ many almost disjoint pieces. This could mean:
    \begin{itemize}
        \item For $X \neq Y$ both subsets of $\lambda$, $S_X \cap S_Y$ is bounded in $\lambda$
        \item For $X \neq Y$ both subsets of $\lambda$, $S_X \cap S_Y$ is non-stationary in $\lambda$.
    \end{itemize}
\end{boxlemma}
\begin{proof}
    For $X \neq Y$, take $\beta = \min(X \symmdiff Y) < \lambda$. Then, for $\alpha > \beta$, $X \cap \alpha \neq Y \cap \alpha$. So we can't have $X \cap \alpha = A_{\alpha}$ and $Y \cap \alpha  = A_{\alpha}$.
\end{proof}

Recall the Solovay Splitting Theorem.

\begin{boxtheorem}[Solovay Splitting Theorem]
    Given $E \ssq \lambda$ stationary, with $\lambda = \cf(\lambda) > \omega$, we can partition $E$ into $\lambda$ many pieces, each of which is stationary.
\end{boxtheorem}
